\subsection{EXAMPLES OF DERIVATIONS}

Example 1. *Let A be an R-algebra of differentiable mappings of R into R and let \(x_0\) be a point of R; R can be considered as an A-module with the external law \((f, a) \mapsto f(x_0)a\) . Then the mapping \(f \mapsto \mathrm{D}f(x_0)\) is a derivation, since (Functions of a Real Variable, I, § 1, no. 3)

\[
(\mathbf {D} (f g)) (x _ {0}) = (\mathbf {D} f (x _ {0})) g (x _ {0}) + f (x _ {0}) (\mathbf {D} g (x _ {0})). *
\]

Example 2. *Let X be a differentiable manifold of class C∞ and let A be the graded R-algebra of differential forms on X. The mapping which associates with every differential form ω on X its exterior differential dω is an anti-derivation of degree +1 (Differentiable and Analytic Manifolds, R, § 8).*

Example 3. Let A be an associative K-algebra. For all \(a \in \mathbf{A}\) , the mapping \(x \mapsto ax - xa\) is a derivation of the algebra A (cf.~no. 6).

Example 4. Let M be a K-module and A the exterior algebra \(\bigwedge (\mathbf{M}^{*})\) with its usual graduation (§ 7, no. 1). *It will be seen in § 11, no. 9 that, for all \(x\in \mathbf{M}\) , the right interior product \(i(x)\) is an antiderivation of A of degree \(-1\).*

Example 5. Returning to the general situation of Definition 1 of no. 2, let \(\overline{\mathbf{K}}\) be another commutative ring and \(\rho: \mathbf{K} \to \overline{\mathbf{K}}\) a ring homomorphism; let \(\overline{\mathbf{A}}, \overline{\mathbf{A}'}\) , \(\overline{\mathbf{A}''}\) , \(\overline{\mathbf{B}}, \overline{\mathbf{B}'}\) , \(\overline{\mathbf{B}''}\) denote the graded \(\overline{\mathbf{K}}\) -modules obtained respectively from \(\mathbf{A}, \mathbf{A}', \mathbf{A}''\) , \(\mathbf{B}, \mathbf{B}'\) , \(\mathbf{B}''\) by extending the ring of scalars to \(\overline{\mathbf{K}}\) (II, § 11, no. 5); we derive from \(\mu\) , \(\lambda_1\) and \(\lambda_2\) \(\overline{\mathbf{K}}\) -bilinear mappings

\[
\bar {\mu}: \overline {{{A}}} \times \overline {{{A}}} ^ {\prime} \rightarrow \overline {{{A}}} ^ {\prime \prime}, \quad \bar {\lambda} _ {1}: \bar {B} \times \overline {{{A}}} ^ {\prime} \rightarrow \bar {B} ^ {\prime \prime}, \quad \bar {\lambda} _ {2}: \overline {{{A}}} \times \bar {B} ^ {\prime} \rightarrow \bar {B} ^ {\prime \prime}
\]

by considering the tensor products by \(l_{\bar{K}}\) of the corresponding K-linear mappings to \(\mu\) , \(\lambda_{1}\) and \(\lambda_{2}\) (II, § 5, no. 1). Then, if d is an \(\varepsilon\) -derivation of degree \(\delta\) of \((\mathbf{A}, \mathbf{A}', \mathbf{A}^{\prime\prime})\) into \((\mathbf{B}, \mathbf{B}', \mathbf{B}^{\prime\prime})\) , the mapping \(\bar{d} = d \otimes l_{\bar{K}}\) of \(\overline{A} \oplus \overline{A}' \oplus \overline{A}''\) into \(\overline{B} \oplus \overline{B}' \oplus \overline{B}''\) is an \(\varepsilon\) -derivation of degree \(\delta\) of \((\overline{A}, \overline{A}', \overline{A}^{\prime\prime})\) into \((\overline{B}, \overline{B}', \overline{B}'' )\) .

Example 6. Let A be a graded K-algebra of type Z; a graded linear K-mapping of degree 0, \(d: A \to A\) , is defined by taking, for \(x_{n} \in \mathbf{A}_{n}(n \in \mathbf{Z})\) , \(d(x_{n}) = nx_{n}\) . This mapping is a derivation of A, since, for \(x_{p} \in A_{p}, x_{q} \in A_{q}\) ,

\[
d (x _ {p} x _ {q}) = (p + q) x _ {p} x _ {q} = d (x _ {p}) x _ {q} + x _ {p} d (x _ {q}).
\]

